% ============================================================================
% DÉFINITION DU COURS 12 : RÉSISTANCE DES MATÉRIAUX
% ============================================================================

\titlehead{Thomas Lusseau - SII en ATS}
\title[RDM]{Résistance des matériaux}
\author[TL]{Thomas Lusseau}
\date{\today}

\TLCourseCoverSetup
  {\repStyle/FondPoly.png}
  {\repStyle/FondPoly.png}
  {Cycle 7 : Dimensionner les pièces d'un mécanisme}
  {Résistance des\\matériaux}
  {Lycée Robert Doisneau - ATS}
  {Thomas Lusseau}

\TLCourseCoverContents{%
  \TLCourseContentsLine{12.1}{Besoin rempli par la RdM}{sec:rdm-1}
  \TLCourseContentsLine{12.2}{Hypothèses de la résistance des matériaux}{sec:rdm-2}
  \TLCourseContentsLine{12.3}{Actions mécaniques extérieures}{sec:rdm-3}
  \TLCourseContentsLine{12.4}{Torseur de cohésion}{sec:rdm-4}
  \TLCourseContentsLine{12.5}{Traction et compression}{sec:rdm-5}
  \TLCourseContentsLine{12.6}{Flexion}{sec:rdm-6}
  \TLCourseContentsLine{12.7}{Sources}{sec:rdm-7}
  \TLCourseContentsLine{12.8}{Exercices du chapitre}{sec:rdm-8}
}

\TLCourseFooter
\TLCourseCover
\markboth{Résistance des matériaux}{Résistance des matériaux}

\livrettrue
\collefalse

\TLStartCourseChapter{12}{Résistance des matériaux}{chap:rdm}

% ============================================================================
% COURS 12 — RÉSISTANCE DES MATÉRIAUX
% Partie A : objectifs, hypothèses, actions extérieures et cohésion
% Source : 16_Rdm_2026.docx
% ============================================================================

\newcommand{\TLRDMCheck}{\ensuremath{\square}}
\newcommand{\TLRDMImage}[2][0.82\linewidth]{%
  \IfFileExists{pandoc/media/#2}{%
    \begin{center}\includegraphics[width=#1]{pandoc/media/#2}\end{center}%
  }{}%
}

\begin{objectifs}
Définir les notions de poutre et les hypothèses fondamentales de la résistance des matériaux. Identifier et caractériser les sollicitations simples. Déterminer un torseur de cohésion, associer un modèle de contraintes à une sollicitation, vérifier la résistance d'une poutre droite et déterminer sa déformée en flexion.
\end{objectifs}

\begin{center}
\small
\begin{tabularx}{\linewidth}{>{\bfseries}p{2.1cm}X c}
\toprule
Je connais & Compétences et connaissances attendues & \\
\midrule
& Les principales sollicitations : traction--compression, cisaillement, torsion, flexion pure et flexion simple. & \TLRDMCheck\\
& L'expression des contraintes et des déformations. & \TLRDMCheck\\
& Le torseur de cohésion. & \TLRDMCheck\\
& Les notions de coefficient de sécurité et de résistance mécanique. & \TLRDMCheck\\
\midrule
Je sais & Poser les hypothèses nécessaires à une étude de RdM : modèle poutre, Saint-Venant, Navier--Bernoulli et linéarité. & \TLRDMCheck\\
& Identifier les contraintes, les déformations et les sollicitations d'un solide. & \TLRDMCheck\\
& Choisir un modèle de solide, indéformable ou déformable, en fonction de l'objectif visé. & \TLRDMCheck\\
& Déterminer le torseur de cohésion dans une poutre. & \TLRDMCheck\\
& Proposer ou justifier les conditions aux limites d'un modèle par éléments finis. & \TLRDMCheck\\
& Déterminer la répartition des contraintes dans une section droite. & \TLRDMCheck\\
& Vérifier la résistance mécanique et le coefficient de sécurité. & \TLRDMCheck\\
& Déterminer l'équation de la flèche d'une poutre droite soumise à un chargement ponctuel ou réparti. & \TLRDMCheck\\
\bottomrule
\end{tabularx}
\end{center}

\section{Besoin rempli par la résistance des matériaux}
\label{sec:rdm-1}

La statique étudie l'équilibre de systèmes de solides supposés indéformables. Cette hypothèse est très restrictive : les pièces réelles se déforment sous l'action des charges et certains problèmes hyperstatiques ne peuvent pas être résolus par la seule statique.

La \textbf{résistance des matériaux}, ou \textbf{RdM}, introduit l'étude des solides déformables lorsqu'ils peuvent être modélisés par des éléments longilignes appelés \textbf{poutres}. Elle établit les relations entre les actions extérieures, les efforts intérieurs, les contraintes, les déformations et les déplacements.

Elle permet notamment :
\begin{itemize}
  \item de calculer ou de vérifier le dimensionnement d'un élément de structure ;
  \item de choisir un matériau à partir de ses caractéristiques mécaniques ;
  \item de contrôler simultanément la résistance et la rigidité ;
  \item d'évaluer l'influence de défauts géométriques ou d'un montage hyperstatique.
\end{itemize}

\begin{center}
\begin{tabular}{ccc}
\IfFileExists{pandoc/media/image5.png}{\includegraphics[width=.29\linewidth]{pandoc/media/image5.png}}{} &
\IfFileExists{pandoc/media/image6.png}{\includegraphics[width=.29\linewidth]{pandoc/media/image6.png}}{} &
\IfFileExists{pandoc/media/image7.png}{\includegraphics[width=.29\linewidth]{pandoc/media/image7.png}}{}\\
\small Géométrie et maillage & \small Déplacements & \small Contraintes
\end{tabular}
\end{center}

\begin{TLWarning}{Deux critères complémentaires}
Une pièce peut ne pas rompre et pourtant être inutilisable si sa déformation est trop importante. Une étude de RdM vérifie donc généralement un critère de \textbf{résistance} et un critère de \textbf{rigidité}.
\end{TLWarning}

\section{Hypothèses de la résistance des matériaux}
\label{sec:rdm-2}

\subsection{Théorie des poutres droites}

Les résultats établis dans ce cours concernent des solides ayant une forme de \textbf{poutre}. Une poutre droite est caractérisée par :
\begin{itemize}
  \item une \textbf{ligne moyenne} continue passant par les barycentres des sections droites ;
  \item une longueur $L$ au moins quatre à cinq fois supérieure à la plus grande dimension $D$ de la section ;
  \item une section constante ou variant progressivement, sans trou ni épaulement brutal dans la zone étudiée ;
  \item un plan de symétrie commun à la géométrie et au chargement dans les problèmes plans.
\end{itemize}

\TLPoutreReperee

\begin{TLWarning}{Poutres constituées de plusieurs tronçons}
Lorsqu'un solide est formé de plusieurs poutres droites, l'étude est décomposée en autant de tronçons rectilignes. Les conditions de continuité relient ensuite les résultats obtenus sur chaque tronçon.
\end{TLWarning}

\subsection{Matériaux de construction}

Dans le modèle de base, le matériau est supposé :
\begin{description}
  \item[homogène :] sa constitution est identique en chaque point ;
  \item[isotrope :] ses propriétés mécaniques sont identiques dans toutes les directions ;
  \item[continu :] les discontinuités microscopiques sont négligées ;
  \item[élastique et linéaire :] après déchargement, la pièce reprend sa forme initiale et les contraintes sont proportionnelles aux déformations.
\end{description}

L'isotropie est une bonne approximation pour de nombreux aciers non fibrés. Elle n'est pas adaptée, sans précaution, au bois, aux matériaux composites ou aux matériaux fortement texturés par laminage ou forgeage.

\subsection{Efforts invariants}

Les déplacements sont supposés petits. Les actions extérieures sont donc exprimées dans le repère lié à la poutre \textbf{non déformée} et sont considérées comme invariantes avant et après application progressive du chargement.

Autrement dit, les calculs sont menés dans la configuration initiale, supposée très proche de la configuration déformée. Cette hypothèse exclut les grandes rotations et les grandes déformations.

\subsection{Hypothèse de Barré de Saint-Venant}

Dans une section située suffisamment loin d'une zone d'application locale des charges, l'état des sollicitations ne dépend que du \textbf{torseur résultant} de ces charges et non de leur répartition détaillée.

Par « suffisamment loin », on entend généralement une distance au moins égale à la plus grande dimension de la section transversale.

\TLRDMSaintVenant

Cette hypothèse permet de remplacer une répartition complexe par un système statiquement équivalent, à condition de ne pas étudier les contraintes au voisinage immédiat de la zone de chargement.

\subsection{Hypothèse de Navier--Bernoulli}

Les sections droites, planes et normales à la ligne moyenne avant chargement, restent planes et normales à la ligne moyenne après déformation. Il n'y a donc pas de gauchissement des sections droites dans le modèle de poutre d'Euler--Bernoulli.

\TLRDMNavierBernoulli

Cette hypothèse convient aux poutres élancées en traction, compression et flexion. Elle devient moins précise pour les poutres courtes, lorsque le cisaillement est important, ou pour la torsion de sections non circulaires.

\subsection{Principe de superposition}

Dans le domaine élastique linéaire, l'effet mécanique d'un ensemble de chargements est égal à la somme des effets produits séparément par chacun d'eux :
\[
\mathcal E(F_1+F_2)=\mathcal E(F_1)+\mathcal E(F_2).
\]

\TLRDMSuperposition

Ce principe peut être appliqué aux réactions d'appui, aux efforts de cohésion, aux contraintes, aux rotations et aux déplacements, tant que les hypothèses de linéarité restent vérifiées.

\section{Modélisation des actions mécaniques extérieures}
\label{sec:rdm-3}

\subsection{Actions mécaniques extérieures}

La modélisation d'un solide par une poutre transforme un problème tridimensionnel en un problème unidimensionnel porté par la ligne moyenne. Les actions mécaniques extérieures sont transposées de la même manière.

Elles peuvent être représentées par :
\begin{itemize}
  \item des \textbf{forces concentrées} appliquées en un point ;
  \item des \textbf{couples concentrés} ;
  \item des \textbf{charges réparties} volumiques, surfaciques ou, après réduction, linéiques ;
  \item les \textbf{actions de liaison} transmises par les appuis et les encastrements.
\end{itemize}

\TLRDMChargesExterieures

\begin{TLWarning}{Répartition des charges}
En RdM, deux systèmes de charges statiquement équivalents ne produisent pas nécessairement le même champ local de contraintes. Le remplacement d'un chargement par un chargement équivalent n'est légitime que dans le domaine d'application de l'hypothèse de Saint-Venant.
\end{TLWarning}

\subsection{Les liaisons}

Les liaisons avec le bâti suppriment certains déplacements et certaines rotations. Elles fournissent à la fois des actions de liaison inconnues et des conditions aux limites utiles à la détermination de la déformée.

\begin{center}
\small
\begin{tabularx}{\linewidth}{c >{\bfseries}l X}
\toprule
Schéma & Liaison & Conséquences dans un problème plan \\
\midrule
\IfFileExists{pandoc/media/image17.png}{\includegraphics[height=9mm]{pandoc/media/image17.png}}{} & Encastrement & bloque les deux translations et la rotation ; transmet deux composantes de force et un moment.\\
\IfFileExists{pandoc/media/image18.png}{
\includegraphics[height=9mm]{pandoc/media/image18.png}}{} & Rotule / appui fixe & bloque les translations ; transmet deux composantes de force.\\
\IfFileExists{pandoc/media/image19.png}{\includegraphics[height=9mm]{pandoc/media/image19.png}}{} & Pivot / articulation & bloque les translations du point de liaison ; laisse libre la rotation.\\
\IfFileExists{pandoc/media/image20.png}{\includegraphics[height=9mm]{pandoc/media/image20.png}}{} & Appui simple & bloque la translation normale à l'appui ; transmet une seule réaction normale.\\
\bottomrule
\end{tabularx}
\end{center}

\section{Actions mécaniques intérieures : torseur de cohésion}
\label{sec:rdm-4}

\subsection{Coupure dans une poutre}

Considérons une poutre $E$ en équilibre sous l'effet d'actions mécaniques extérieures. Une \textbf{coupure imaginaire}, réalisée dans un plan perpendiculaire à la ligne moyenne, sépare la poutre en deux tronçons :
\[
E=E_1\cup E_2.
\]

Les actions exercées par $E_2$ sur $E_1$ à travers la section droite $S$ sont des actions mécaniques intérieures. Leur résultante est représentée, au barycentre $G$ de la section, par le \textbf{torseur de cohésion} :
\[
\left\{\mathcal T_{\mathrm{coh}}\right\}_G
=\left\{\mathcal T_{E_2\rightarrow E_1}\right\}_G.
\]

\TLCoupurePoutre

\subsection{Détermination du torseur de cohésion}

Deux isolements sont possibles.

\paragraph{Isolement du tronçon gauche $E_1$.}
Le principe fondamental de la statique donne :
\[
\left\{\mathcal T_{\mathrm{ext}\rightarrow E_1}\right\}_G
+\left\{\mathcal T_{E_2\rightarrow E_1}\right\}_G=\{0\},
\]
d'où :
\[
\boxed{\left\{\mathcal T_{\mathrm{coh}}\right\}_G
=-\left\{\mathcal T_{\mathrm{ext}\rightarrow E_1}\right\}_G.}
\]

\paragraph{Isolement du tronçon droit $E_2$.}
Par actions réciproques et équilibre de $E_2$ :
\[
\boxed{\left\{\mathcal T_{\mathrm{coh}}\right\}_G
=\left\{\mathcal T_{\mathrm{ext}\rightarrow E_2}\right\}_G.}
\]

Les notations $E_1/E_2$, gauche/droite ou $x^-/x^+$ sont équivalentes. Le torseur est toujours réduit au barycentre $G$ de la section considérée.

\begin{TLMethod}{Déterminer le torseur de cohésion}
\begin{enumerate}
  \item Déterminer les réactions de liaison par l'équilibre global de la poutre.
  \item Définir les intervalles d'étude et effectuer une coupure à l'abscisse $x$.
  \item Isoler le tronçon qui contient le moins de charges.
  \item Transporter toutes les actions extérieures au point $G$ de la coupure.
  \item Appliquer le PFS et identifier les six composantes du torseur de cohésion.
  \item Vérifier la continuité des diagrammes et les sauts produits par les charges concentrées.
\end{enumerate}
\end{TLMethod}

Dans la base locale $(\vec x,\vec y,\vec z)$ liée à la poutre :
\[
\left\{\mathcal T_{\mathrm{coh}}\right\}_G=
\left\{
\begin{array}{c}
N\vec x+T_y\vec y+T_z\vec z\\[1mm]
M_t\vec x+M_{fy}\vec y+M_{fz}\vec z
\end{array}
\right\}_G.
\]

\begin{center}
\small
\begin{tabularx}{\linewidth}{>{\bfseries}c X}
\toprule
Composante & Signification \\
\midrule
$N$ & effort normal, perpendiculaire à la section droite ; $N>0$ en traction et $N<0$ en compression.\\
$T_y$, $T_z$ & efforts tranchants, contenus dans le plan de la section.\\
$M_t$ & moment de torsion, porté par l'axe de la poutre.\\
$M_{fy}$, $M_{fz}$ & moments fléchissants autour des axes principaux de la section.\\
\bottomrule
\end{tabularx}
\end{center}

\subsection{Nature des sollicitations}

Une \textbf{sollicitation simple} correspond à un torseur de cohésion ne comportant qu'une composante principale. Une \textbf{sollicitation composée} associe plusieurs sollicitations simples.

\begin{center}
\small
\begin{tabularx}{\linewidth}{>{\bfseries}l c c c c X}
\toprule
Sollicitation & $N$ & $T$ & $M_t$ & $M_f$ & Commentaire \\
\midrule
Traction / compression & $\neq0$ & 0 & 0 & 0 & $N>0$ en traction, $N<0$ en compression.\\
Cisaillement simple & 0 & $\neq0$ & 0 & 0 & modèle idéal où le moment parasite est négligé.\\
Torsion simple & 0 & 0 & $\neq0$ & 0 & rotation des sections autour de la ligne moyenne.\\
Flexion pure & 0 & 0 & 0 & $\neq0$ & moment fléchissant constant, effort tranchant nul.\\
Flexion simple & 0 & $\neq0$ & 0 & $\neq0$ & cas courant des poutres chargées transversalement.\\
Sollicitation composée & variable & variable & variable & variable & combinaison de plusieurs composantes.\\
\bottomrule
\end{tabularx}
\end{center}

\subsection{Vision locale : contrainte en un point}

Le torseur de cohésion donne une vision globale des actions transmises par une section. La vision locale décrit leur répartition sur chaque élément de surface $\mathrm dS$.

Au point $M$ d'une section orientée par la normale extérieure $\vec n$, le \textbf{vecteur contrainte} est défini par :
\[
\vec t(M,\vec n)=\lim_{\Delta S\to0}\frac{\Delta\vec f}{\Delta S}.
\]

Il se décompose en une composante normale et une composante tangentielle :
\[
\vec t=\sigma\vec n+\vec\tau,
\qquad
\sigma=\vec t\cdot\vec n,
\qquad
\vec\tau=\vec t-(\vec t\cdot\vec n)\vec n.
\]

\TLRDMVecteurContrainte

Une contrainte s'exprime en pascals. En pratique :
\[
\boxed{1\ \mathrm{MPa}=1\ \mathrm{N/mm^2}.}
\]

% ============================================================================
% COURS 12 — RÉSISTANCE DES MATÉRIAUX
% Partie B : traction, flexion et déformée
% Source : 16_Rdm_2026.docx
% ============================================================================

\section{Traction et compression}
\label{sec:rdm-5}

\subsection{Définition}

Une poutre est sollicitée en \textbf{traction simple} lorsque les actions appliquées à ses extrémités se réduisent à deux forces égales, opposées et portées par la ligne moyenne. L'effort intérieur associé est l'\textbf{effort normal} $N$.

\TLDiagrammeTraction

Dans toute section droite :
\[
\left\{\mathcal T_{\mathrm{coh}}\right\}_G
=\left\{\begin{array}{c}N\vec x\\\vec0\end{array}\right\}_G.
\]

Par convention :
\[
N>0\quad\text{en traction},
\qquad
N<0\quad\text{en compression}.
\]

\subsection{Contrainte normale}

Sous l'hypothèse d'une répartition uniforme sur une section droite d'aire $S$ :
\[
N=\iint_S\sigma\,\mathrm dS,
\qquad
\boxed{\sigma=\frac{N}{S}.}
\]

La contrainte est positive en traction et négative en compression. Dans une zone comportant un accident géométrique, la contrainte maximale est évaluée avec un coefficient de concentration de contrainte $K_t$ :
\[
\sigma_{\max}=K_t\frac{|N|}{S_{\mathrm{nette}}}.
\]

\subsection{Condition de résistance}

On note :
\begin{itemize}
  \item $R_e$ la limite d'élasticité du matériau ;
  \item $s>1$ le coefficient de sécurité ;
  \item $R_{pe}=R_e/s$ la résistance pratique à l'extension ;
  \item $K_t$ le coefficient de concentration de contrainte.
\end{itemize}

La condition de résistance s'écrit :
\[
\boxed{\sigma_{\max}=K_t\frac{|N|}{S_{\mathrm{nette}}}
\leq R_{pe}=\frac{R_e}{s}.}
\]

Le coefficient de sécurité effectif est :
\[
s_{\mathrm{eff}}=\frac{R_e}{\sigma_{\max}}.
\]

\subsection{Déformations longitudinales et transversales}

Pour une poutre de longueur initiale $L_0$ et de longueur déformée $L$ :
\[
\Delta L=L-L_0,
\qquad
\boxed{\varepsilon_L=\frac{\Delta L}{L_0}.}
\]

Le coefficient de Poisson $\nu$ relie la déformation transversale à la déformation longitudinale :
\[
\boxed{\varepsilon_T=-\nu\varepsilon_L.}
\]

Pour l'acier, on utilise fréquemment $\nu\simeq0{,}3$.

\subsection{Loi de Hooke}

Dans le domaine élastique linéaire :
\[
\boxed{\sigma=E\varepsilon_L,}
\]
où $E$ est le module d'Young. On en déduit, pour une poutre prismatique soumise à un effort normal constant :
\[
\boxed{\Delta L=\frac{NL_0}{ES}.}
\]

Pour l'acier, $E$ est voisin de $210\,000\ \mathrm{MPa}$.

\begin{center}
\IfFileExists{pandoc/media/image43.png}{\includegraphics[width=.48\linewidth]{pandoc/media/image43.png}}{}
\end{center}

\begin{TLWarning}{Domaine de validité}
La loi de Hooke ne décrit que le domaine élastique linéaire. Au-delà de la limite d'élasticité, la déformation devient en partie permanente et la superposition n'est plus applicable.
\end{TLWarning}

\subsection{Application guidée : tirant}

\begin{TLMethod}{Dimensionner un tirant en acier}
Un tirant de longueur $L=2\ \mathrm m$ supporte un effort normal de traction $N=5\,000\ \mathrm N$. Le matériau possède $R_e=300\ \mathrm{MPa}$ et $E=200\,000\ \mathrm{MPa}$. On retient $s=1{,}7$.

\begin{enumerate}
  \item La section minimale vérifie $N/S\leq R_e/s$, donc :
  \[
  S_{\min}=\frac{sN}{R_e}.
  \]
  \item Pour une section circulaire :
  \[
  d_{\min}=\sqrt{\frac{4S_{\min}}{\pi}}.
  \]
  \item L'allongement est ensuite calculé avec :
  \[
  \Delta L=\frac{NL}{ES}.
  \]
\end{enumerate}
\end{TLMethod}

\section{Flexion}
\label{sec:rdm-6}

\subsection{Définition et cinématique}

Une poutre est soumise à de la \textbf{flexion pure} lorsque le torseur de cohésion se réduit à un moment fléchissant. En flexion simple, un effort tranchant accompagne le moment fléchissant.

Sous l'hypothèse de Navier--Bernoulli, une section droite reste plane et normale à la ligne moyenne. Les fibres longitudinales situées d'un côté de la fibre neutre s'allongent, tandis que celles situées de l'autre côté se raccourcissent.

Pour une courbure $1/\rho$ et une fibre située à l'ordonnée $y$ :
\[
\varepsilon_{xx}(y)=-\frac{y}{\rho}.
\]

Avec la loi de Hooke :
\[
\sigma_{xx}(y)=E\varepsilon_{xx}(y)=-E\frac{y}{\rho}.
\]

\TLDistributionContrainteFlexion

\subsection{Contrainte normale de flexion}

Le moment fléchissant est le moment de la répartition des contraintes normales :
\[
M_{fz}=\iint_S -y\sigma_{xx}(y)\,\mathrm dS.
\]

En introduisant le moment quadratique :
\[
I_{Gz}=\iint_S y^2\,\mathrm dS,
\]
on obtient :
\[
\boxed{\sigma_{xx}(y)=-\frac{M_{fz}}{I_{Gz}}y.}
\]

La relation moment--courbure est :
\[
\boxed{M_{fz}=\frac{E I_{Gz}}{\rho}.}
\]

La contrainte maximale est atteinte sur la fibre la plus éloignée de l'axe neutre :
\[
|\sigma_{\max}|=\frac{|M_{fz}|}{I_{Gz}}y_{\max}.
\]

En introduisant le module de flexion $W_z=I_{Gz}/y_{\max}$ :
\[
\boxed{|\sigma_{\max}|=\frac{|M_{fz}|}{W_z}.}
\]

\subsection{Condition de résistance}

Avec un coefficient de concentration de contrainte $K_f$ :
\[
\boxed{K_f\frac{|M_{fz}|}{W_z}\leq R_{pe}=\frac{R_e}{s}.}
\]

Lorsque la fonction impose également une limite de déplacement, on vérifie :
\[
\boxed{|f_{\max}|\leq f_{\mathrm{limite}}.}
\]

\begin{TLWarning}{Résistance et rigidité}
Le dimensionnement en flexion est gouverné soit par la contrainte maximale, soit par la flèche maximale. Le critère le plus contraignant dépend de la géométrie, du matériau et du cahier des charges.
\end{TLWarning}

\subsection{Équation différentielle de la déformée}

Pour de petites rotations, la courbure est assimilée à la dérivée seconde de la déformée $y(x)$ :
\[
\frac{1}{\rho}\simeq y''(x).
\]

Avec la convention de signe retenue dans le problème :
\[
\boxed{E I_{Gz}\,y''(x)=M_{fz}(x).}
\]

L'intégration conduit successivement à la rotation et à la flèche :
\[
y'(x)=\int\frac{M_{fz}(x)}{EI_{Gz}}\,\mathrm dx+C_1,
\]
\[
y(x)=\iint\frac{M_{fz}(x)}{EI_{Gz}}\,\mathrm dx^2+C_1x+C_2.
\]

Les constantes sont déterminées à partir des conditions aux limites :
\begin{itemize}
  \item encastrement en $x=0$ : $y(0)=0$ et $y'(0)=0$ ;
  \item appui simple en $x=a$ : $y(a)=0$ ;
  \item continuité en $x=a$ : $y_1(a)=y_2(a)$ et $y'_1(a)=y'_2(a)$.
\end{itemize}

\TLDiagrammeFlexion

\subsection{Résultats classiques de déformée}

Les formules ci-dessous supposent $E$ et $I_{Gz}$ constants.

\begin{center}
\small
\renewcommand{\arraystretch}{1.55}
\begin{tabularx}{\linewidth}{>{\bfseries}p{3.4cm} X X}
\toprule
Chargement & Moment fléchissant maximal & Flèche maximale \\
\midrule
Poutre sur deux appuis, charge ponctuelle $F$ au milieu &
$\displaystyle |M_{\max}|=\frac{FL}{4}$ &
$\displaystyle |f_{\max}|=\frac{FL^3}{48EI_{Gz}}$ \\
Poutre sur deux appuis, charge répartie uniforme $q$ &
$\displaystyle |M_{\max}|=\frac{qL^2}{8}$ &
$\displaystyle |f_{\max}|=\frac{5qL^4}{384EI_{Gz}}$ \\
Poutre encastrée, charge ponctuelle $F$ à l'extrémité &
$\displaystyle |M_{\max}|=FL$ &
$\displaystyle |f_{\max}|=\frac{FL^3}{3EI_{Gz}}$ \\
Poutre encastrée, charge répartie uniforme $q$ &
$\displaystyle |M_{\max}|=\frac{qL^2}{2}$ &
$\displaystyle |f_{\max}|=\frac{qL^4}{8EI_{Gz}}$ \\
\bottomrule
\end{tabularx}
\end{center}

\begin{center}
\begin{tabular}{cc}
\IfFileExists{pandoc/media/image47.png}{\includegraphics[width=.42\linewidth]{pandoc/media/image47.png}}{} &
\IfFileExists{pandoc/media/image50.png}{\includegraphics[width=.42\linewidth]{pandoc/media/image50.png}}{}\\[2mm]
\IfFileExists{pandoc/media/image53.png}{\includegraphics[width=.42\linewidth]{pandoc/media/image53.png}}{} &
\IfFileExists{pandoc/media/image56.png}{\includegraphics[width=.42\linewidth]{pandoc/media/image56.png}}{}
\end{tabular}
\end{center}

\begin{TLMethod}{Résoudre un problème de flexion}
\begin{enumerate}
  \item Modéliser la poutre, les appuis et les charges.
  \item Déterminer les réactions de liaison par le PFS.
  \item Calculer $T_y(x)$ et $M_{fz}(x)$ sur chaque intervalle.
  \item Tracer les diagrammes et localiser $|M_{fz}|_{\max}$.
  \item Vérifier la contrainte maximale avec $\sigma=M_{fz}y/I_{Gz}$.
  \item Intégrer $EI_{Gz}y''=M_{fz}$ et appliquer les conditions aux limites.
  \item Vérifier le critère de flèche du cahier des charges.
\end{enumerate}
\end{TLMethod}

\section{Sources}
\label{sec:rdm-7}

Ce cours a été élaboré à partir de documents industriels publics et d'activités pédagogiques de collègues de l'UPSTI. Les exercices de synthèse conservent les références indiquées dans le document d'origine, notamment Jérôme Letard et le concours commun PT-SI 2011.

% ============================================================================
% COURS 12 — EXERCICES, PARTIE A
% Source : 16_Rdm_2026.docx
% ============================================================================

\section{Exercices du chapitre}
\label{sec:rdm-8}

\begin{center}
\IfFileExists{pandoc/media/image59.png}{\includegraphics[width=.52\linewidth]{pandoc/media/image59.png}}{}
\end{center}

\subsection{Traction -- compression}

\subsubsection{Exercice 1 : choix d'un câble}

\begin{center}
\begin{tabular}{cc}
\IfFileExists{pandoc/media/image61.png}{
\includegraphics[width=.43\linewidth]{pandoc/media/image61.png}}{} &
\IfFileExists{pandoc/media/image62.png}{\includegraphics[width=.43\linewidth]{pandoc/media/image62.png}}{}
\end{tabular}
\end{center}

On souhaite soulever une poutre de béton armé de longueur $6\ \mathrm m$ et de masse $2\ \mathrm t$. Deux sangles sont disposées symétriquement et reliées à une élingue constituée de deux brins de longueur $l=4\ \mathrm m$.

Chaque câble est constitué de six torons de dix-neuf fils de diamètre $0{,}8\ \mathrm{mm}$, entourant une âme textile dont la résistance mécanique est négligée. La limite d'élasticité d'un fil est $R_e=1770\ \mathrm{N/mm^2}$.

\begin{enumerate}
  \item Déterminer la tension maximale dans un brin de câble en fonction de la géométrie de l'élingue.
  \item Le coefficient de sécurité imposé aux structures de levage étant $s=6$, déterminer le diamètre minimal du câble à choisir.
  \item Déterminer le coefficient de sécurité effectif du câble retenu.
\end{enumerate}

\subsubsection{Exercice 2 : câble de puits de mine}

\begin{minipage}{.72\linewidth}
On cherche la forme d'égale résistance d'un câble soumis à son propre poids et supportant à son extrémité une nacelle d'effort $F=2\times10^3\ \mathrm{daN}$.

Données :
\begin{itemize}
  \item poids volumique : $\omega=7{,}8\ \mathrm{daN/dm^3}$ ;
  \item résistance pratique : $R_{pe}=100\ \mathrm{MPa}$ ;
  \item longueur : $L=500\ \mathrm m$.
\end{itemize}
\end{minipage}\hfill
\begin{minipage}{.24\linewidth}
\IfFileExists{pandoc/media/image63.png}{\includegraphics[width=\linewidth]{pandoc/media/image63.png}}{}
\end{minipage}

Pour un câble d'égale résistance, la section varie suivant :
\[
S(x)=\frac{F}{R_{pe}}\exp\!\left(\frac{\omega x}{R_{pe}}\right),
\]
avec des unités cohérentes.

\begin{enumerate}
  \item Déterminer les diamètres aux deux extrémités du câble d'égale résistance.
  \item Déterminer le diamètre d'un câble de section constante supportant la même nacelle.
  \item Calculer et comparer les masses des deux câbles.
\end{enumerate}

\subsubsection{Exercice 3 : bielle de moteur Diesel et chemise mince}

\paragraph{Compression de la bielle.}
Une bielle de moteur Diesel est soumise à un effort maximal de compression de $150\times10^3\ \mathrm N$. Sa partie centrale est modélisée par un prisme de section $S=300\ \mathrm{mm^2}$ et de longueur $L=180\ \mathrm{mm}$.

\begin{enumerate}
  \item Avec un coefficient de sécurité $s=1{,}2$, déterminer la limite d'élasticité minimale de l'acier spécial utilisé.
  \item Le module d'Young étant $E=2\times10^5\ \mathrm{MPa}$, calculer le raccourcissement maximal de la bielle.
\end{enumerate}

\paragraph{Paroi mince de la chemise.}
Le piston a un diamètre intérieur $d=155\ \mathrm{mm}$ et la pression effective dans la chambre vaut $p=8\ \mathrm{MPa}$. La chemise est assimilée à un tube de longueur $L$, de diamètre intérieur $d$ et d'épaisseur $e$ faible devant $d$. La fonte utilisée possède $R_{pe}=50\ \mathrm{MPa}$.

La coupure diamétrale conduit à un effort normal :
\[
N=pLd,
\]
et la section résistante vaut approximativement $2eL$.

\begin{enumerate}[resume]
  \item Déterminer l'épaisseur minimale de la chemise.
\end{enumerate}

\subsubsection{Exercice 4 : maillon de chaîne}

\begin{minipage}{.66\linewidth}
La plaquette d'une chaîne à rouleaux est soumise à la traction. Le perçage du plan $A-A$ produit un coefficient de concentration $K_{tp}=1{,}6$ ; la réduction de section du plan $B-B$ est associée à $K_{tr}=1{,}05$.

Caractéristiques de l'acier :
\[
R_m=1300\ \mathrm{MPa},\qquad
R_e=1100\ \mathrm{MPa},\qquad
E=2{,}2\times10^5\ \mathrm{MPa}.
\]
\end{minipage}\hfill
\begin{minipage}{.30\linewidth}
\IfFileExists{pandoc/media/image65.png}{\includegraphics[width=\linewidth]{pandoc/media/image65.png}}{}
\end{minipage}

\begin{enumerate}
  \item Évaluer la charge limite de rupture statique $F_{\mathrm{lim}}$ en recherchant la section la plus critique.
  \item Pour un fonctionnement dynamique avec $s=3$, déterminer la charge maximale de service $F_{\max}$.
\end{enumerate}

\subsection{Fardeleuse}

\textit{Source : concours commun PT-SI 2011.}

\begin{center}
\begin{tabular}{cc}
\IfFileExists{pandoc/media/image66.jpeg}{\includegraphics[width=.36\linewidth]{pandoc/media/image66.jpeg}}{} &
\IfFileExists{pandoc/media/image67.png}{
\includegraphics[width=.52\linewidth]{pandoc/media/image67.png}}{}
\end{tabular}
\end{center}

Le fardelage consiste à déposer un film plastique thermo-rétractable autour d'un produit puis à le chauffer. Le produit est déplacé par un tapis entraîné par adhérence entre un cylindre moteur et un cylindre tendeur.

La tension initiale $T_0$ est réglée en translatant l'axe du cylindre tendeur. Le cylindre est guidé par deux roulements qui exercent, en $B$ et $C$, deux efforts :
\[
\vec F=-F\vec y,
\qquad F=230\ \mathrm N.
\]
L'axe est modélisé par une poutre cylindrique de diamètre $d$, reposant sur deux appuis en $A$ et $D$.

\begin{center}
\begin{tabular}{cc}
\IfFileExists{pandoc/media/image68.png}{\includegraphics[width=.38\linewidth]{pandoc/media/image68.png}}{} &
\IfFileExists{pandoc/media/image69.png}{\includegraphics[width=.50\linewidth]{pandoc/media/image69.png}}{}
\end{tabular}
\end{center}

\subsubsection{Dimensionnement de l'axe}

\begin{enumerate}
  \item Établir la relation entre la tension $T_0$ du tapis et la norme $F$ des efforts en $B$ et $C$.
  \item Déterminer les réactions des appuis en $A$ et $D$ en fonction de $F$.
  \item Exprimer les composantes du torseur de cohésion le long de la poutre et tracer les diagrammes correspondants. Identifier la sollicitation.
  \item En négligeant l'effort tranchant, localiser la zone la plus sollicitée.
  \item Déterminer la contrainte normale maximale $\sigma_{xx,\max}$.
  \item Avec un coefficient de sécurité $s_c=2$, établir l'expression du diamètre minimal $d_{\min}$.
\end{enumerate}

\subsubsection{Choix du matériau}

Les barres disponibles ont un diamètre compris entre $3$ et $50\ \mathrm{mm}$.

\begin{center}
\small
\begin{tabular}{lllrr}
\toprule
Acier & État & & $R_e$ (MPa) & $R_m$ (MPa)\\
\midrule
C22 & normalisé && 240 & 430\\
C22 & trempé et revenu && 340 & 650\\
C55 & normalisé && 370 & 680\\
C55 & trempé et revenu && 550 & 950\\
15 Cr Ni 6 & trempé et revenu && 650 & 1000\\
25 Cr Mo 4 & trempé et revenu && 550 & 850\\
\bottomrule
\end{tabular}
\end{center}

\begin{enumerate}[resume]
  \item Choisir un acier et un état métallurgique compatibles avec le diamètre disponible et le coefficient de sécurité.
  \item Pour l'acier choisi, calculer $d_{\min}$ puis sélectionner le diamètre commercial immédiatement supérieur.
\end{enumerate}

% ============================================================================
% COURS 12 — EXERCICES, PARTIE B
% Questions de cours, autocontraintes et QCM
% Source : 16_Rdm_2026.docx
% ============================================================================

\subsection{Questions de cours}

\begin{enumerate}
  \item Préciser la signification de l'acronyme RdM et citer les deux objectifs principaux d'une étude de résistance des matériaux.
  \item Citer les cinq principales hypothèses qui permettent de simplifier une étude de poutre.
  \item Écrire le torseur de cohésion complet et donner le nom de ses six composantes.
  \item Donner la définition du vecteur contrainte $\vec t(M,\vec n)$.
  \item Écrire la condition de résistance d'une pièce soumise à la traction et préciser tous les termes utilisés.
  \item Citer six sollicitations simples ou usuelles.
  \item Écrire les lois de Hooke en contrainte normale et en cisaillement. Donner le nom des modules $E$ et $G$ ainsi que leur relation pour un matériau isotrope :
  \[
  G=\frac{E}{2(1+\nu)}.
  \]
  \item Compléter le vocabulaire associé aux symboles suivants : $\nu$, $\theta$, $\varepsilon$, $\rho$, $\sigma$, $\tau$ et $\gamma$.
  \item Associer chaque sollicitation à la déformation et au torseur de cohésion correspondants.
  \item Dans un schéma de chargement, préciser le nom de l'action $F$ et son unité.
  \item Écrire l'expression de la contrainte normale en traction puis en flexion.
  \item Donner le nom et l'unité de $I_{Gz}$, $y$, $M_{fz}$ et $R_{pe}$.
  \item Donner l'équation différentielle de la déformée en flexion.
  \item Donner les conditions aux limites usuelles :
  \begin{itemize}
    \item encastrement en $x=0$ ;
    \item appui simple en $x=a$ ;
    \item continuité entre deux tronçons au point $x=a$.
  \end{itemize}
\end{enumerate}

\subsection{Défauts géométriques et autocontraintes}

L'arbre 2 est lié au carter 1 par deux roulements. En $O_1$, un roulement à deux rangées de billes à contacts obliques est modélisé par une liaison pivot d'axe $\vec x_1$. En $A$, un roulement à billes à contact radial dont la bague extérieure est libre axialement est modélisé par une liaison linéaire annulaire.

Après fabrication, les deux alésages du carter présentent un défaut de coaxialité radial $\delta_r$. Le montage force l'arbre à se déformer et crée des actions de liaison sans qu'aucune charge extérieure ne soit appliquée : ce sont des \textbf{autocontraintes}.

\begin{center}
\IfFileExists{pandoc/media/image80.png}{\includegraphics[width=.70\linewidth]{pandoc/media/image80.png}}{}
\end{center}

On suppose le problème plan. Le plan est choisi de sorte que la composante hors plan du défaut soit nulle. Les composantes des torseurs de liaison sont notées $X_O$, $Y_O$, $N_O$ au point $O_1$ et $Y_A$ au point $A$.

\subsubsection{Étude statique et hyperstatisme}

\begin{enumerate}
  \item Écrire, dans le problème plan, les torseurs des actions de liaison du carter sur l'arbre aux points $O_1$ et $A$.
  \item Faire le bilan des inconnues et déterminer le degré d'hyperstatisme interne de la liaison équivalente entre l'arbre et le carter.
  \item Aucune charge extérieure n'agissant sur l'arbre, appliquer le PFS et établir les trois relations entre les inconnues de liaison.
\end{enumerate}

\subsubsection{Déformée de l'arbre}

La section de l'arbre possède un moment quadratique $I_{Gz}$ et le matériau un module d'Young $E$.

\begin{enumerate}[resume]
  \item Montrer que le moment fléchissant, pour $0\leq x\leq L$, peut s'écrire :
  \[
  M_{fz}(x)=-N_O+xY_O.
  \]
  \item En utilisant les conditions d'encastrement équivalent au point $O_1$, déterminer l'expression de la déformée $y(x)$.
  \item Utiliser la condition géométrique $y(L)=\delta_r$ pour lever l'hyperstaticité et montrer que :
  \[
  \boxed{Y_A=\frac{3EI_{Gz}\delta_r}{L^3}.}
  \]
  Déterminer ensuite $Y_O$ et $N_O$.
  \item En déduire l'expression complète de $y(x)$ en fonction de $E$, $I_{Gz}$, $L$ et $\delta_r$.
\end{enumerate}

\subsubsection{Application numérique}

On donne :
\[
D=20\ \mathrm{mm},\qquad L=300\ \mathrm{mm},\qquad
E=2\times10^5\ \mathrm{MPa},\qquad \delta_r=0{,}05\ \mathrm{mm}.
\]

\begin{enumerate}[resume]
  \item Calculer le moment quadratique :
  \[
  I_{Gz}=\frac{\pi D^4}{64}.
  \]
  \item Calculer les actions de liaison.
  \item Déterminer le moment fléchissant maximal puis l'autocontrainte normale maximale dans l'arbre.
\end{enumerate}

\subsection{QCM de synthèse}

Une seule réponse est attendue pour chaque question. Pour les applications numériques, on pourra prendre $\pi=3$ et $g=10\ \mathrm{m/s^2}$ lorsque cela simplifie les calculs.

\begin{enumerate}
  \item Comment s'appelle le diagramme représentant la contrainte en fonction de la déformation lors d'un essai de traction ?
  \begin{itemize}
    \item loi de Hooke ; essai de traction ; diagramme contrainte--déformation ; module d'Young.
  \end{itemize}

  \item Sur une courbe de traction conventionnelle, comment nomme-t-on la zone comprise entre la limite d'élasticité et la rupture ?
  \begin{itemize}
    \item domaine élastique ; domaine plastique ; striction ; écrouissage.
  \end{itemize}

  \item Parmi les affirmations suivantes, laquelle caractérise correctement un acier dont la limite d'élasticité vaut $295\ \mathrm{MPa}$ ?
  \begin{itemize}
    \item au-delà de cette contrainte, les déformations peuvent devenir permanentes ;
    \item la contrainte est alors nécessairement tangentielle ;
    \item la loi $\sigma=E\varepsilon$ reste valable jusqu'à la rupture.
  \end{itemize}

  \item Quel est l'ordre de grandeur du module d'Young du plexiglas ?
  \begin{itemize}
    \item $3{,}1\ \mathrm{GPa}$ ; $62\ \mathrm{GPa}$ ; $210\ \mathrm{GPa}$.
  \end{itemize}

  \item Déterminer la coordonnée du centre d'aire de la section composée représentée ci-dessous.
  \begin{center}
  \IfFileExists{pandoc/media/image82.png}{\includegraphics[width=.30\linewidth]{pandoc/media/image82.png}}{}
  \end{center}

  \item Quelle affirmation concernant le torseur de cohésion est fausse ?
  \begin{itemize}
    \item $M_f$ désigne un moment fléchissant ;
    \item $T_z$ est l'effort tranchant suivant $\vec z$ ;
    \item $N$ est l'effort tranchant suivant $\vec y$ ;
    \item $M_t$ désigne le moment de torsion.
  \end{itemize}

  \item Une poutre dont le torseur de cohésion ne comporte que $M_{fz}$ est soumise à :
  \begin{itemize}
    \item traction ; compression ; cisaillement ; torsion ; flexion pure ; flexion simple.
  \end{itemize}

  \item Quel terme n'est pas une grandeur ou une notion usuelle de RdM ?
  \begin{itemize}
    \item résistance mécanique ; contrainte tangentielle ; moment fléchissant ; module de Coulomb ; pression uniforme.
  \end{itemize}

  \item Dans un torseur de cohésion, que représente $N$ ?
  \begin{itemize}
    \item un effort normal ; un moment nominal ; une unité en newtons.
  \end{itemize}

  \item Identifier la sollicitation principale de la pièce représentée.
  \begin{center}
  \IfFileExists{pandoc/media/image84.png}{
\includegraphics[width=.48\linewidth]{pandoc/media/image84.png}}{}
  \end{center}

  \item Associer chacune des déformations usuelles au torseur de cohésion réduit correspondant : traction, cisaillement, torsion, flexion pure et flexion simple.

  \item En quelle unité s'exprime un moment quadratique ou un moment polaire de section ?
  \begin{itemize}
    \item $\mathrm{mm^2}$ ; $\mathrm{mm^3}$ ; $\mathrm{mm^4}$ ; $\mathrm{N/mm^2}$.
  \end{itemize}

  \item Une éprouvette reçoit une force normale de $52\ \mathrm{kN}$ sur une section de $260\ \mathrm{mm^2}$. Son module d'Young vaut $40\ \mathrm{GPa}$. Calculer la déformation relative.
  \begin{itemize}
    \item $0{,}3\ \%$ ; $0{,}5\ \%$ ; $1{,}5\ \%$ ; $2\ \%$.
  \end{itemize}

  \item Quelle lettre désigne le premier coefficient de Lamé ?
  \begin{itemize}
    \item $\nu$ ; $\varepsilon$ ; $\lambda$ ; $\theta$ ; $\rho$ ; $\tau$ ; $\gamma$ ; $\sigma$.
  \end{itemize}

  \item Dans une section droite soumise à de la flexion pure, le torseur de cohésion est :
  \begin{itemize}
    \item nul ; un glisseur ; un couple ; un torseur quelconque.
  \end{itemize}

  \item Une force normale de $48\ \mathrm N$ est uniformément répartie sur une section circulaire de rayon $12\ \mathrm{cm}$. Calculer la contrainte moyenne.
  \begin{itemize}
    \item $120\ \mathrm{Pa}$ ; $1{,}1\ \mathrm{kPa}$ ; $6{,}7\ \mathrm{kPa}$ ; $280\ \mathrm{kPa}$.
  \end{itemize}

  \item Une poutre horizontale de longueur $L$, encastrée à une extrémité, reçoit une force verticale $F$ à l'autre extrémité. Quel est le moment fléchissant maximal ?
  \begin{itemize}
    \item $FL$ ; $FL/2$ ; $FL/4$ ; $FL^2/2$.
  \end{itemize}

  \item Une section circulaire de diamètre $100\ \mathrm{mm}$ reçoit un moment fléchissant de $75\ \mathrm{N\,m}$. Calculer la contrainte normale maximale.
  \begin{itemize}
    \item $160\ \mathrm{Pa}$ ; $400\ \mathrm{kPa}$ ; $0{,}8\ \mathrm{MPa}$ ; $12{,}8\ \mathrm{MPa}$.
  \end{itemize}

  \item Comment modélise-t-on le poids propre d'une poutre droite homogène ?
  \begin{itemize}
    \item par un glisseur unique au centre de gravité ;
    \item par une densité linéique de force ;
    \item par des moments appliqués aux appuis.
  \end{itemize}

  \item Parmi les propositions suivantes, laquelle correspond à une sollicitation composée ?
  \begin{itemize}
    \item extension ; cisaillement ; flexion déviée ; torsion simple.
  \end{itemize}

  \item Dans un essai idéal de cisaillement suivant $\vec y$, quelle composante parasite du torseur de cohésion doit être négligée pour parler de cisaillement simple ?
  \begin{itemize}
    \item $N$ ; $T_y$ ; $T_z$ ; $M_t$ ; $M_{fy}$ ; $M_{fz}$.
  \end{itemize}

  \item Une poutre circulaire de rayon $12\ \mathrm{mm}$ reçoit un moment de torsion de $9\ \mathrm{N\,m}$. Calculer la contrainte tangentielle maximale avec :
  \[
  \tau_{\max}=\frac{M_tR}{J},
  \qquad
  J=\frac{\pi R^4}{2}.
  \]
  \begin{itemize}
    \item $3{,}5\ \mathrm{kPa}$ ; $3{,}5\ \mathrm{MPa}$ ; $6{,}9\ \mathrm{MPa}$ ; $27{,}8\ \mathrm{MPa}$.
  \end{itemize}
\end{enumerate}


\printindex
